% Pursuing Stacks (cotes 134-2 to 134-8): archive cote 134-4, site page 32 (PDF page 33), typescript page 120.
% typescript. Draft compared: ps3.tex:1584-1642.
% AI-transcribed by Claude Opus 5.5 (claude-opus-5-5, adaptive thinking, effort high), Message Batches API, batch ps2 (msgbatch_01NVwr1Yuo8qJf5rykeGBXPK), 2026-11-12; not yet checked by eye.
\scanpage{134-4/32}
\mspage{120}
there should not exist such a generating interval,  the natural next best
assumption to make is the existence of a generating set $\hand{\Sigma}_h^{\circ}$ which is
``small''  (namely \inserted{an} element of the``universe''we are working \xout{with} in ).  The
case of a finite generating set of intervals reduces to the case of a
single one though, \deleted{at least when M is stable under products,} by just
taking the product of those intervals.

An interesting case is when the generating set $\hand{\Sigma}_h^{\circ}$ consists of
contractible objects of M.  Such a generating set exists iff the structure
considered $\hand{\Sigma}_h$ comes from a contractibility structure.  About the best we
could hope for is the existence of a \inserted{single} \uline{generating} \uline{contractible} \xout{element in}
\uline{weak homotopy interval} $\underline{I}$ .  If we got any interval $\underline{I}$ in M,  this can be
viewed as a generating contractible weak homotopy \xout{\unsure{structure on}} interval
for a suitable homotopy structure on M (then necessarily unique) iff
\xout{\unsure{I satisfies th}} the identity \inserted{map} \xout{\unsure{functor}} of I can be joined to a constant
one by a chain of maps, such that two consecutive ones are \inserted{tied by a} $\underline{I}$-homotopy\xout{\unsure{x}} \inserted{or by a $\check{\underline{I}}$-homotopy.)}
The most evident way to meet this condition is by a \uline{one-step} \uline{chain} from
$\mathrm{id}_I$ to the constant map defined by one of the endpoints, $\hand{\delta}_1$ say.  This
brings us back to the structure of a composition law
\[ IxI \hand{\longrightarrow} I \]
in I,  having $e_o$ as a left unit and $e_1$ as a left zero\xout{\unsure{x}} element.  Let's
\margin{X} call an interval,  endowed with such a composition law,  a\deleted{n} \uline{\over{\unsure{interactor}}{contractor}}
in M.  Thus starting from a\deleted{n} \over{\unsure{interactor}}{contractor} in M is about the nicest way to
define a homotopy structure in M,  as a matter of fact the strongest type
of such a structure  - namely a contractibility structure,  admitting a
\uuline{generating} \uuline{contractible} \uuline{weak homotopy interval}  (and better still, admit-%
ting a \uline{generating} \uline{contractor}). \margin{//} Of course,  starting with the weakest
kind of homotopy structure on M,  namely just a homotopy relation
$R_h \in \mathrm{Hom}_1(M)$,  if $\underline{I}$ is a homotopy interval which is moreover endowed with
\xout{the} a structure of a contractor,  \xout{the} i.e. \inserted{if $\underline{I}$ is} a contra\unsure{d|c}tor such that the
\inserted{end-point} sections $\hand{\delta}_o\ ,\hand{\delta}_1$  are homotopic,  then I is automatically contractible
(never mind if it is generating or not).  \xout{\unsure{I get the imp}}

It seems to me that the homotopy structures I've looked at so far
(such as (Cat) and various standart elementary modelizers A\^{}) are not only
contractibility structures,  but they \xout{can} all can be defined by a single
contractor each.

Besides the \hand{``}\uline{basic} \uline{contractor}\hand{''}$\hand{\Delta}_1$ in (Cat),  there are two general ways
I've met so far for getting contractors.  \unsure{One} \xout{is} has been made explicit in
these notes a number of times,  namely the \uline{Lawvere element} \xout{in} $L_M$  in M

\margin{32}
